\begin{abstract}
Many tokenised assets, including funds, structured products, and corporate bonds, are valued only by their issuer, its administrator, or a primary dealer, so their value is a periodic published valuation rather than a price quoted every block. Decentralised lending's margining assumptions fit neither that valuation basis nor a liquidation that often takes weeks. Where a secondary market exists, liquidity in large amounts is rarely available within a block, and the only guaranteed exit is redemption with the issuer. D-SIMM applies the International Swaps and Derivatives Association's Standard Initial Margin Model (SIMM) to margining in decentralised finance (DeFi). Off-chain risk is what SIMM was built to margin. In this first iteration we focus on the most popular case, tokenised funds. We derive narrow bounds from published risk weights for loan-to-value (LTV) limits, liquidation thresholds (LLTV, liquidation loan-to-value), and the borrow and supply caps. A liquidated position is assumed to be held through the clearing window to cash settlement. Free parameters remain for the practitioner to choose, but the aim is to reduce the impact of discretion on leverage in DeFi as far as possible, by setting transparent limits.
\end{abstract}
